\documentclass{article}
\usepackage{array}
\usepackage{caption}
\usepackage{graphicx}
\usepackage{siunitx}
\usepackage{tabularx}
\usepackage{rotating}
\usepackage{booktabs}
\usepackage{threeparttable}
\sisetup{output-exponent-marker = e,round-mode = figures, round-precision = 3,
	scientific-notation = true}
\newcolumntype{L}[1]{>{\hsize=#1\hsize\raggedright\arraybackslash}X}%
\newcolumntype{R}[1]{>{\hsize=#1\hsize\raggedleft\arraybackslash}X}%
\newcolumntype{C}[1]{>{\hsize=#1\hsize\centering\arraybackslash}X}%
\begin{document}
	\begin{threeparttable}
		\captionsetup{labelformat=empty}
		\caption{Summary of ANOVA test}
		\setlength{\tabcolsep}{4pt}
		\begin{tabularx}{\textwidth}{L{1} R{1} R{1} R{1} R{1}}
			
			\toprule
			
			Model Comparison & Difference (km) & Lower Bound (km) & Upper Bound (km) & Adj pvalue \\
			
			\midrule
			
			62-46 $z_{1100}$ & 8.25 & 2.52 & 14.0 & 2e-03 \\
			78-46 $z_{1100}$ & 18.0 & 12.3 & 23.7 & 1e-10 \\
			94-46 $z_{1100}$ & 33.6 & 27.8 & 39.3 & 2e-11 \\
			78-62 $z_{1100}$ & 9.75 & 4.02 & 15.5 & 2e-04 \\
			94-62 $z_{1100}$ & 25.3 & 19.6 & 31.0 & 2e-11 \\
			94-78 $z_{1100}$ & 15.6 & 9.84 & 21.3 & 7e-09 \\
			
			\bottomrule

\end{tabularx}
\begin{tablenotes}
\item Tukey's post-hoc test for comparing the means of coupling depth between groups
of models with different backarc lithospheric thickness. Low pvalues indicate means
are actually different
\end{tablenotes}
\end{threeparttable}
\end{document}